\begin{afterword}
\summaryhead

The post argues that the idea of religion as a ``separate magisterium,'' beyond proof or disproof, is a modern invention. Ancient religion made factual claims and staged tests of them. Elijah's contest with the priests of Baal is retold as an experiment, and the revelation at Sinai as unmistakable evidence if it happened. The Old Testament, written with no libraries to check it, made large historical claims such as the Exodus, which archaeology has not found. The New Testament, written under Roman record-keeping, claims smaller miracles, but still offers them as evidence.

Religion once ruled law, history, sexual morals, government and science, and ethics is what remained when better institutions took the rest. But the Bible's ethics, the author argues, are as obsolete as its science, and people endorse them only because they can get away with it. The post concludes that the idea of religion as non-disprovable is a ``Big Lie,'' false to history, to scripture, to what children are taught and to what most believers hold.

\responsehead

The post's historical point is sound and well chosen. Ancient scriptures do present their claims as public events, and some, like the Exodus, have been checked against the record and found wanting. The trouble is what the post aims that point at, and how far it stretches it.

The target is never named. The phrase ``separate magisterium'' comes from Stephen Jay Gould's essay of 1997, which argued that science and religion should be treated as having separate domains of authority, and claimed that this was already the standard attitude of the major Western religions. Most of the post's evidence answers neither claim. Its evidence from history, scripture and what children are told shows what religions have taught, not how the domains should be related. Its fourth item, what most believers on Earth hold, does meet a claim of Gould's, but the post gives no evidence for it. The courtroom joke at the end assumes the bloody axe is in hand. For the Exodus it is; for Gould's own example, the soul, there is none. Calling the idea a ``Big Lie'' puts ridicule where an answer to Gould's argument was needed.

The section on ethics changes the subject. That the Bible approves of slavery is a moral judgment, not an observation. It shows that religious ethics can be criticized, not that they can be disproved by evidence, which is the post's thesis. The claim that ``Most people's concept of rationality is determined by what they think they can get away with'' is offered with no evidence at all.

The scope outruns the examples. Every religious example comes from the Bible, yet the claims cover ``The vast majority of religions in human history'' and what is ``strictly Western.'' Some details are also loose. Bartley's ``retreat to commitment'' is given without the argument it named, that every position rests on some arbitrary starting point. The New Testament does report big public wonders, such as the darkness at the crucifixion and the risen saints appearing in Jerusalem. And the footnote misstates the Old Testament twice: Deuteronomy 22:5 sets no death penalty for cross-dressing, and the Psalms and Job contain passages of wonder at the heavens and at nature.

In short: a sound point about ancient religion, turned against an unnamed position whose argued form it does not answer, with claims about all religions drawn from the Bible alone.
\end{afterword}
